\documentclass[11pt]{article}
\usepackage[margin=1in]{geometry}
\usepackage{amsmath,amssymb,amsthm}
\usepackage[hidelinks]{hyperref}
\newtheorem{theorem}{Theorem}
\newcommand{\ket}[1]{\lvert #1\rangle}
\newcommand{\bra}[1]{\langle #1\rvert}
\newcommand{\norm}[1]{\lVert #1\rVert}
\newcommand{\BPP}{\mathsf{BPP}}
\newcommand{\BQP}{\mathsf{BQP}}
\title{A fixed Pauli-norm threshold for diagonal matrix powering}
\author{Daniel Fredriksen\\[2pt]\small Quantyra Inc., Las Vegas, Nevada, USA\\
\small \href{https://orcid.org/0009-0002-9123-7435}{ORCID 0009-0002-9123-7435}}
\date{}
\begin{document}
\maketitle
\begin{abstract}
Given a Hermitian contraction as an explicit sum of Pauli strings, consider estimating a single diagonal entry of its polynomial power to inverse-polynomial additive error. Cifuentes et al.\ asked whether this problem remains BQP-complete at constant Pauli coefficient norm. We show that the threshold promise problem is PromiseBQP-complete for every fixed norm cap \(c>1\), even for positive definite contractions with dyadic coefficients. At cap one an efficient classical sampler is known. The reduction mixes a cyclic circuit clock with the identity. This preserves an eigenvalue whose spectral weight encodes acceptance, while the rest of the spectrum can still be suppressed by a polynomial power. We give a finite-bit reduction that preserves the two norm promises and its decision gap.
\end{abstract}

\section{The fixed-cap question}
For a Hermitian qubit operator \(B=\sum_P b_PP\), let \(\lambda(B)=\sum_P|b_P|\), where \(P\) ranges over distinct tensor products of \(I,X,Y,Z\), including the identity. Write \(\norm{B}\) for the operator norm. Janzing and Wocjan encoded a quantum circuit decision into a diagonal entry of a power of a sparse matrix~\cite{jw}. Cifuentes et al.\ gave a Pauli description of such a clock with polynomial coefficient norm and asked whether a constant norm suffices~\cite[Discussion and Sec.~IV\,A]{cif}. Their constant induced matrix \(1\)-norm bound concerns column sums in the computational basis. Here the bound is on coefficients in the Pauli basis.

Fix \(c>0\). The input lists real dyadic coefficients and their Pauli labels, a computational basis index \(j\), a positive integer \(m\), a rational threshold \(g\), and a positive tolerance \(\epsilon\). Both \(m\) and \(1/\epsilon\) are polynomially bounded in the input length, with a fixed polynomial bound or unary resource limits. The promises are \(\norm{B}\le1\), \(\lambda(B)\le c\), and one of
\[
   \bra j B^m\ket j\ge g+\epsilon,\qquad
   \bra j B^m\ket j\le g-\epsilon .
\]
The decision is YES in the first case. Comparing an estimate of additive error \(\epsilon/2\) with \(g\) decides the promised input.

\begin{theorem}\label{thm:main}
For \(0<c\le1\) the diagonal-power threshold problem has a bounded-error classical polynomial-time algorithm. For every fixed \(c>1\) it is PromiseBQP-complete under deterministic classical polynomial-time many-one reductions. The hard instances can be chosen positive definite with dyadic coefficients and an explicitly specified basis input.
\end{theorem}

The result concerns fixed caps. In particular, the statement distinguishes cap one from every fixed cap above one under the same inverse-polynomial accuracy convention. It does not classify every cap approaching one as a function of input length. We give the classical endpoint and quantum upper bound for completeness, followed by the construction establishing hardness.

\section{Two estimation algorithms}
For \(\lambda(B)\le1\), assign probability \(|b_P|\) to each Pauli label and place any remaining mass on a dummy label. These dyadic probabilities can be sampled exactly using their common power-of-two denominator. Draw \(m\) independent labels. Score zero on any dummy draw. Otherwise score the real part of the basis diagonal of the ordered product of the sampled Paulis times the product of their coefficient signs. Pauli multiplication and basis-diagonal evaluation take polynomial time. Expanding the expectation gives \(\bra j B^m\ket j\). The sample belongs to \([-1,1]\), so \(O(\epsilon^{-2})\) repetitions suffice for additive error \(\epsilon/2\) with success at least \(2/3\). This is the Pauli-sampling bound used in~\cite{cif}.

For the quantum upper bound, use the spectral-sampling algorithm of~\cite{jw,cif}. Prepare \(\ket j\) and perform phase estimation for \(e^{iB}\). Its eigenphases in signed radians lie in \([-1,1]\). Estimate an eigenvalue \(x\) to \(\epsilon/(8m)\) with failure at most \(\epsilon/32\), clip to \([-1,1]\), and raise the result to \(m\). Conditional on an eigenstate, the usual confidence amplification of phase estimation applies; on \(\ket j\) it samples the corresponding spectral mixture. Since \(x\mapsto x^m\) is \(m\)-Lipschitz on this interval, the bias is at most \(\epsilon/8+2\epsilon/32\) before scalar rounding. With an additional \(\epsilon/16\) rounding allowance, \(O(\epsilon^{-2})\) independent trials still estimate the diagonal to error less than \(\epsilon\) with success at least \(2/3\). Run at half tolerance for the decision problem.

All required simulation times are polynomial in \(m/\epsilon\). For example, a first-order product formula applies each listed Pauli rotation in each of \(r\) slices. Its operator error at time \(t'\) is at most \(t'^2\lambda(B)^2/(2r)\) by the sum-of-commutators bound. Controlled Pauli rotations and their finite-precision synthesis thus have polynomial gate cost. Allocate total simulation and synthesis error at most \(\epsilon/128\) across a trial, and use ideal phase-estimation failure at most \(\epsilon/64\); these fit the stated failure allowance. Evaluating the scalar power to the prescribed accuracy takes polynomial bit time.

\section{A clock with a distinguished eigenvalue}
Start with a bounded-error quantum circuit \(C=G_t\cdots G_1\) over \(H,T,\mathrm{CNOT}\) and their inverses, on a work register initially in a basis state \(\ket z\). Let \(\Pi\) accept output bit one and let \(p=\bra z C^\dagger\Pi C\ket z\). Put \(M=2^{\lceil\log_2\max(4,2t+1)\rceil}\). On a cyclic clock with \(M\) positions, take in order the gates of \(C\), the reflection \(2\Pi-I=-Z_{\rm out}\), the gates of \(C^\dagger\), and identity padding. Call this list \(V_0,\ldots,V_{M-1}\). Its ordered product is \(R=C^\dagger(2\Pi-I)C\). On the binary clock and work register set
\[
 W=\sum_{a=0}^{M-1}\ket{a+1\bmod M}\bra a\otimes V_a,\qquad
 A=\tfrac12(W+W^\dagger),\qquad \ket j=\ket0\ket z .
\]
Every clock label has one predecessor and one successor, and each \(V_a\) is unitary. Hence \(W\) is unitary on the full space and \(A\) is a Hermitian contraction.

If \(Rv=\sigma v\), \(\sigma\in\{+1,-1\}\), the vectors \(W^a\ket{0,v}\) for \(0\le a<M\) form an orthonormal ring for unit \(v\). The last edge of the ring carries the sign \(\sigma\). Its Fourier modes have eigenphases \(2\pi k/M\) for \(\sigma=+1\) and \((2k+1)\pi/M\) for \(\sigma=-1\), with \(0\le k<M\). Each mode has squared overlap \(1/M\) with \(\ket{0,v}\). An orthonormal eigenbasis of \(R\) gives orthogonal rings spanning the whole space. Also \((I+R)/2=C^\dagger\Pi C\), so the positive sector has weight \(p\) at the initial state. Taking real parts of the clock eigenvalues shows that the spectral measure of \(A\) at \(j\) has mass exactly \(p/M\) at \(+1\). Every other eigenvalue is at most
\[
 \cos(\pi/M)\le 1-2/M^2 ,
\]
using \(1-\cos u=2\sin^2(u/2)\). This is the cyclic-clock spectral calculation of~\cite{jw}, with the chosen reflection assigning acceptance to the periodic sector.

\section{Enforcing the cap and retaining the gap}
A clock matrix unit is a tensor product of one-bit units. Each one-bit unit has two Pauli coefficients of magnitude \(1/2\); each clock transition therefore has \(M\) Pauli terms of total weight one. Every gate used above has Pauli coefficient sum at most two: \(H=(X+Z)/\sqrt 2\), \(T\) is diagonal with coefficients \((1\pm e^{i\pi/4})/2\), and \(\mathrm{CNOT}\) has four coefficients of magnitude \(1/2\). The same bounds hold for adjoints. Multiplying on disjoint registers, summing over transitions and taking the Hermitian part gives an explicit list for \(A\) with \(O(M^2)\) terms and \(\lambda(A)\le K=2M\).

Choose a fixed inverse power of two \(d\le\min(1/2,c-1)\), and define
\[
 s=d/K,\qquad B_*=(1-s)I+sA.
\]
Here \(0<s\le1/4\), \((1-2s)I\le B_*\le I\), and
\[
 \lambda(B_*)\le1-s+sK = 1-s+d\le c.
\]
The affine map fixes \(+1\). All eigenvalues other than the distinguished value lie in \([0,1-2s/M^2]\). Choose the integer
\[
 m=\frac{M^2}{2s}\,\lceil\log_2(64M)\rceil=O(M^3\log M/d).
\]
The contribution from the complement of \(+1\) is nonnegative and at most \(\exp(-2sm/M^2)\le1/(64M)\). Thus
\[
 \frac{p}{M}\le \bra j B_*^m\ket j\le \frac{p}{M}+\frac{1}{64M}.
\]

\paragraph{Finite coefficients.}
Set \(\tau=2^{-\lceil\log_2(128mM)\rceil}\). Multiply the collected list for \(B_*\) by \(1-\tau\), then round each real coefficient toward zero on a dyadic grid with spacing at most \(\tau/(2S)\), where \(S\ge1\) is the number of listed coefficient slots. Call the result \(B\). The coefficients of the Hermitian part before rounding belong to \(\mathbb Z[1/2]+\sqrt{2}\,\mathbb Z[1/2]\). Only one constant-size gate enters each expanded clock term, and combining terms uses polynomially many additions. Directed rounding on this number field takes polynomial bit time by integer sign comparisons and squaring. The grid requires \(O(\log(SmM))\) fractional bits.

Rounding reduces the absolute coefficient sum. Its operator error relative to \((1-\tau)B_*\) is at most \(\tau/2\), by summing the individual errors. Consequently \(\lambda(B)\le c\), \(\norm{B}\le1-\tau/2\), and its least eigenvalue is at least \((1-\tau)(1-2s)-\tau/2>0\). Telescoping the powers of the two contractions gives
\[
 \norm{B^m-B_*^m}\le m\norm{B-B_*}\le \tfrac32m\tau\le\frac{3}{256M}.
\]
Use the threshold \(g=1/(2M)\) and margin \(\epsilon=1/(16M)\). On YES instances \(p\ge2/3\) and on NO instances \(p\le1/3\). Even allowing acceptance error \(1/48\) from compiling a different standard source gate set, the scaled diagonal \(M\bra j B^m\ket j\) is greater than \(7/12\) on YES and less than \(5/12\) on NO. A single estimate to \(1/(32M)\) therefore decides the source input.

All Pauli labels and coefficients are emitted in classical polynomial time. For every fixed \(c>1\), both \(m\) and the inverse margin are polynomial. If the resource bounds are stated in the number of qubits, one may append \(M\) idle work qubits with identity labels throughout. This leaves the operator norm, coefficient sum and measured diagonal unchanged. This completes the many-one reduction and the proof of Theorem~\ref{thm:main}.

\paragraph{Acknowledgment.} The author used OpenAI Codex to assist with derivation, drafting and reproducibility checks.
\begin{thebibliography}{2}
\bibitem{jw} D. Janzing and P. Wocjan.
A simple PromiseBQP-complete matrix problem.
\emph{Theory of Computing} \textbf{3}, 61--79 (2007).
\href{https://doi.org/10.4086/toc.2007.v003a004}{doi:10.4086/toc.2007.v003a004}.
\bibitem{cif} S. Cifuentes, S. Wang, T. L. Silva, M. Berta and L. Aolita.
Quantum computational complexity of matrix functions.
\emph{PRX Quantum} \textbf{7}, 020364 (2026).
\href{https://doi.org/10.1103/g5x4-jcsz}{doi:10.1103/g5x4-jcsz}.
\end{thebibliography}
\end{document}
